\documentclass[11pt,a4paper]{article}
\usepackage[a4paper,margin=25mm,headheight=15pt]{geometry}
\usepackage{fontspec}
\defaultfontfeatures{Ligatures=TeX}
% Load TeX Live fonts by filename: arXiv has limited fontconfig registration.
\setmainfont{lmroman10-regular.otf}[
BoldFont=lmroman10-bold.otf,ItalicFont=lmroman10-italic.otf,
BoldItalicFont=lmroman10-bolditalic.otf]
\setmonofont[Scale=0.88]{lmmono10-regular.otf}
\usepackage{amsmath,unicode-math,microtype,enumitem,titlesec,fancyhdr,xurl,booktabs}
\setmathfont{latinmodern-math.otf}
\usepackage[hidelinks,unicode]{hyperref}
\hypersetup{pdftitle={ROS3-Limassol: An OpenPELAOS Prototype for Telemetry Liveness Monitoring},
pdfauthor={Charalampos Nadiotis},
pdfsubject={Technical report; initial idea credited to Christos Paparistodemou}}
\linespread{1.02}
\setlength{\parindent}{0pt}
\setlength{\parskip}{4pt}
\setlength{\emergencystretch}{2em}
\setlist{nosep,leftmargin=1.4em}
\titleformat{\section}{\large\bfseries}{\thesection}{0.65em}{}
\titleformat{\subsection}{\normalsize\bfseries}{\thesubsection}{0.65em}{}
\titlespacing*{\section}{0pt}{12pt}{5pt}
\titlespacing*{\subsection}{0pt}{8pt}{3pt}
\pagestyle{fancy}
\fancyhf{}
\fancyhead[L]{\small OpenPELAOS / ROS3-Limassol}
\fancyhead[R]{\small Update: 31 October 2026}
\fancyfoot[C]{\small\thepage}
\renewcommand{\headrulewidth}{0.3pt}
\newcommand{\code}[1]{\texttt{#1}}
\newcommand{\healthy}{\mathrm{HEALTHY}}
\newcommand{\stale}{\mathrm{STALE}}
\newcommand{\waiting}{\mathrm{WAITING}}
\begin{document}
\thispagestyle{plain}
\begin{center}
{\LARGE\bfseries ROS3-Limassol: An OpenPELAOS Prototype\\[3pt]
for Telemetry Liveness Monitoring}\\[8pt]
{\large Intermittently Connected Robotic Systems}\\[11pt]
Charalampos Nadiotis\\[3pt]
OpenPELAOS, independent open-source project\\[5pt]
{\small Initial project idea: Christos Paparistodemou}\\[5pt]
{\small 23 October 2026 \quad|\quad Preliminary technical report}
\end{center}

\begin{abstract}
Robotic telemetry can stop arriving because of communication disruption, publisher
failure, or upstream data interruption. A receiver therefore needs an explicit
rule for interpreting the absence of valid messages. This report describes
ROS3-Limassol, an OpenPELAOS development prototype comprising a transport-independent
Python core, a synthetic unmanned surface vehicle telemetry source, and an optional
ROS~2 adapter. The observer evaluates elapsed receiver time since the last accepted
message and distinguishes waiting, healthy, and stale states. Schema validation,
session identifiers, and increasing sequence numbers prevent several malformed
or replayed messages from refreshing that timer. A deterministic experiment
with a three-second blackout reproduces stale detection and recovery, accepting
51 of 81 generated samples. In the original version-0.1.0 validation, eighteen core tests pass; one ROS graph test is
explicitly skipped because ROS~2 is unavailable in the authoring environment.
The evidence establishes software behavior under a specified synthetic scenario.
Real transport timing, field performance, and research novelty remain unestablished.
An additional dated section summarizes subsequent synthetic and ArduPilot SITL
evidence; the reported SITL aggregates await an independent raw-recording audit
and make no physical-link or field claim.
\end{abstract}
{\small\textbf{Keywords:} ROS~2; field robotics; telemetry; liveness monitoring;
intermittent connectivity; reproducibility.}

\textbf{Scope of this review draft (31 October 2026).}
Sections 1--7 preserve the original version-0.1.0 account, including its
then-pending work. Section~\ref{sec:simulation-update} and
Appendix~\ref{app:simulation-tables} describe later simulation evidence
separately. This draft is for author review, not a submission or an assertion
that the historical experiments have been independently reproduced.

\section{Introduction}
A robotics operator may continue to see an apparently valid position or battery
reading after the source has stopped delivering usable updates. The most recent
value and the current availability of its source are different observations.
ROS3-Limassol begins with the task of making arrival liveness explicit, before
attempting navigation, mission execution, or broader autonomy.

The project is developed under OpenPELAOS, with the initial idea credited to
Christos Paparistodemou. ROS3-Limassol is the project identifier; version~0.1.0
implements a toolkit that can use ROS~2. The immediate outcome is a small
implementation that can be inspected and exercised without a robot or a ROS
installation. Its intended ROS mode adds interprocess communication around the
same core logic. This is an independent project, rather than an official
ROS distribution.

The motivating question is whether a bounded receiver-side observer can report
the absence and resumption of accepted telemetry without depending on synchronized
sender and receiver clocks. The present contribution is an implementation and
deterministic verification baseline \cite{artifact}. It does not introduce a new
failure-detection algorithm or establish an advantage over existing robotics
software. Its value at this stage is a precise contract, reproducibility, and
an explicit boundary between verified behavior and planned evaluation.

\section{Background and technical basis}
ROS~2 provides an established modular robotics software architecture.
Macenski et al.\ describe its design and deployment context
\cite{macenski2022}. Reusing that infrastructure allows the prototype to focus on
application-level observations. The adapter follows the Python publisher/subscriber
and \code{ament\_python} conventions in the official documentation
\cite{ros_python}.

ROS~2 quality-of-service settings distinguish history, reliability, and durability,
and communicating endpoints require compatible policies \cite{ros_qos}.
These policies do not, by themselves, define what a particular application should
consider a sufficiently recent valid telemetry sample. The prototype therefore
evaluates its own acceptance and timeout rules.

Failure-detector theory distinguishes completeness and accuracy, including the
possibility that a detector makes mistakes \cite{chandra1996}. That distinction
helps interpret a timeout as a receiver observation of missing accepted data.
It cannot identify which physical component failed, and no formal
failure-detector class is claimed here.

Receiver intervals use Python's \code{time.monotonic()}. Its documented reference
point is unspecified, while differences between readings are meaningful and
unaffected by system-clock updates \cite{python_time}. Sender timestamps remain
recording metadata. Local arrival liveness is therefore separate from a future
model of measurement age or clock synchronization.

\section{System model and liveness rule}
\subsection{Accepted arrivals}
Consider one configured robot and a sequence of accepted receiver arrival times
$r_0,r_1,\ldots$ on the same monotonic clock. The sequence is non-decreasing.
For observation time $t$, define
\begin{equation}
\mathcal{R}(t)=\{r_i:r_i\leq t\}.
\label{eq:arrivals}
\end{equation}
This set contains arrivals that have passed the application checks. In the
default scenario, the last accepted arrival before the blackout is
$r_i=1.9\,\mathrm{s}$. A malformed payload received later does not enter this set.

When $\mathcal{R}(t)$ is nonempty, the most recent accepted arrival has age
\begin{equation}
a(t)=t-\max\mathcal{R}(t).
\label{eq:age}
\end{equation}
The value is an interval in seconds, rather than the age of the underlying
sensor measurement. At $t=2.9\,\mathrm{s}$, with a last arrival at
$1.9\,\mathrm{s}$, Equation~\eqref{eq:age} gives $a(t)=1.0\,\mathrm{s}$.

For timeout $\tau>0$, the observer state is
\begin{equation}
s(t)=
\begin{cases}
\waiting,&\mathcal{R}(t)=\varnothing,\\
\healthy,&\mathcal{R}(t)\neq\varnothing\ \text{and}\ a(t)<\tau,\\
\stale,&\mathcal{R}(t)\neq\varnothing\ \text{and}\ a(t)\geq\tau.
\end{cases}
\label{eq:state}
\end{equation}
With $\tau=1.0\,\mathrm{s}$, the previous example is stale. A new accepted message
at $5.0\,\mathrm{s}$ resets the age to zero and restores the healthy state.
Before the first valid arrival, the observer retains the waiting state.

\subsection{Observation timing}
Suppose no accepted arrival occurs after $r_\ast$, and consecutive observer
evaluations are separated by at most $h$ seconds. The first stale report satisfies
\begin{equation}
r_\ast+\tau\leq t_{\mathrm{detect}}\leq r_\ast+\tau+h.
\label{eq:bound}
\end{equation}
For $r_\ast=1.9\,\mathrm{s}$, $\tau=1.0\,\mathrm{s}$, and an assumed evaluation-gap
bound $h=0.1\,\mathrm{s}$, the interval is $[2.9,3.0]\,\mathrm{s}$. The deterministic
scenario detects the transition at the lower endpoint.

Equation~\eqref{eq:bound} depends on the evaluation-gap assumption. A nominal ROS
timer period does not establish a hard scheduling guarantee on a general operating
system. No real-time bound has been measured. Furthermore, a stream of old
measurements carrying new valid sequence numbers can satisfy arrival liveness
while remaining unsuitable for navigation.

\section{Prototype architecture}
\subsection{Telemetry contract}
The core accepts a versioned JSON object containing robot and publisher-session
identifiers, sequence number, timestamp, geodetic latitude and longitude, heading,
speed, and battery percentage. Units are explicit: degrees, metres per second,
percentage, and timestamp nanoseconds. The source produces illustrative signals,
rather than a physically consistent vessel dynamics model.

RFC~8259 specifies JSON, excludes non-finite numeric values, and discusses
interoperability concerns including duplicate object names \cite{rfc8259}.
The application contract adds strict field and value checks. All expected fields
must be present; extra fields, duplicate keys, and booleans used as numbers are
rejected. Numeric ranges and identifier lengths are bounded.

Payloads are limited to 4096 UTF-8 bytes before JSON decoding. The version-1 speed
limit is $100\,\mathrm{m\,s^{-1}}$, which is an application bound. These checks
bound application parsing and retained state; they do not establish middleware
resource bounds.

\subsection{Sessions and ordering}
Within the active publisher session, accepted sequence numbers must strictly
increase. Gaps are permitted. A previously unseen session permits a publisher
restart with sequence zero. Known retired sessions are rejected, preventing
delayed samples from a known old session from becoming current again.

The monitor retains at most 128 sessions by default. An additional unseen session
is rejected once the limit is reached, while the active session may continue.
Restarting the monitor clears the history. Persistent replay protection and
authenticated session negotiation remain future work.

The model assumes one legitimate publisher for the configured robot identity.
It does not authenticate sources or resolve competing publishers. Invalid
payloads, unexpected robots, duplicate or out-of-order sequences, retired sessions,
and excess new sessions increment a rejection counter without refreshing liveness.

\subsection{Transport separation}
The validator, simulator, monitor, and local demonstration use only Python's
standard library. The default demonstration passes JSON directly between source
and observer within one process.

The optional adapter carries that JSON in \code{std\_msgs/msg/String}.
Telemetry endpoints use KEEP\_LAST, depth~1, BEST\_EFFORT, and VOLATILE.
Status uses depth~1, RELIABLE, and VOLATILE. These settings follow documented
ROS~2 semantics \cite{ros_qos}; compatible transport profiles do not establish
measurement freshness.

In ROS mode, the observer is configured to evaluate and publish status every
$0.1\,\mathrm{s}$ using a steady-clock timer, with separate robot namespaces.
A status consumer still needs to detect disappearance of the observer itself.
Version~0.1.0 contains no vehicle-command interface, generated ROS message
definitions, or MAVLink adapter.

\section{Verification method and results}
\subsection{Deterministic experiment}
The scenario uses duration $D=8.0\,\mathrm{s}$, source rate $f=10\,\mathrm{Hz}$,
and receiver timeout $1.0\,\mathrm{s}$. Sampling includes both endpoints:
\begin{equation}
t_k=\frac{k}{f},\qquad k=0,\ldots,\lfloor Df\rfloor.
\label{eq:samples}
\end{equation}
The configured values give indices $0,\ldots,80$, hence 81 generated samples.

The link model suppresses samples for which $2.0\leq t_k<5.0\,\mathrm{s}$.
These are indices 20 through 49, so
\begin{equation}
N_{\mathrm{generated}}=81,\quad
N_{\mathrm{suppressed}}=30,\quad
N_{\mathrm{accepted}}=81-30=51.
\label{eq:counts}
\end{equation}
Suppression occurs before delivery to the observer, explaining why the base
scenario has zero rejected messages despite missing 30 generated samples.

The observed sequence is waiting at $0.0\,\mathrm{s}$ before the first sample,
healthy at $0.0\,\mathrm{s}$ after acceptance, stale at $2.9\,\mathrm{s}$, and healthy
again at $5.0\,\mathrm{s}$ \cite{artifact}. The final pre-blackout sequence is 19
and the first resumed sequence is 50.

Detection is $2.9-2.0=0.9\,\mathrm{s}$ after the configured blackout starts, or
$2.9-1.9=1.0\,\mathrm{s}$ after the last accepted arrival. These intervals use
different reference events. Recovery at the $5.0\,\mathrm{s}$ virtual tick
demonstrates the state rule rather than a measured network recovery latency.

\subsection{Executed software checks}
The recorded run discovers 19 test methods: 18 core methods pass and one ROS graph
test is skipped because \code{rclpy} is unavailable. The core checks cover timeout
boundaries, recovery, malformed payloads, invalid clocks, sequence ordering,
publisher restarts, retired sessions, history bounds, JSON validation, and
deterministic command-line output \cite{artifact}. This is evidence for the
exercised behaviors rather than a complete correctness proof or coverage result.

An offline wheel build succeeds. Inspection confirms the core modules, ROS
adapter, launch file, package manifest, license, and ament marker. The installed
wheel's JSON demonstration also runs from outside the source tree. Python syntax,
the XML package manifest, and CI YAML parse successfully.

The ROS package build, launch execution, DDS graph test, and ROS CI job were not
executed in the authoring environment. The supplied graph test observes healthy
traffic, cancels the publisher timer to model an outage, and resets it to check
recovery when run in a suitable ROS environment. The test definition alone does
not establish that the transport has passed.

\section{Limitations and next evaluation}
The deterministic scenario uses virtual time and in-process delivery. It does not
exercise discovery delay, scheduling, serialization overhead, queueing, bandwidth
constraints, or a physical wireless link. Arrival age and state labels do not
certify measurement capture time, position accuracy, or vehicle condition.

The next evaluation should execute the Jazzy build and graph test, then record
a real ROS launch with the same blackout. Controlled experiments can vary delay,
jitter, outages, and loss while documenting source rates, session handling, and
timeouts. Detection and recovery delays require explicitly defined ground-truth
events; stale indications during source availability should also be reported.
Repeated runs are needed before estimating performance distributions.

Recorded-data replay, a receive-only MAVLink adapter, and companion-computer
deployment are proposed extensions. Multi-robot field operation, marine dynamics,
new transport performance, and middleware replacement have not been demonstrated.
Research novelty requires a broader literature comparison and additional evidence.

\section{Conclusion}
ROS3-Limassol establishes an OpenPELAOS starting point for explicit telemetry
arrival-liveness monitoring. A receiver monotonic clock, strict message acceptance,
and bounded session history make its behavior inspectable. The deterministic
blackout reproduces stale detection and recovery, and 18 core test methods pass.
A successful ROS transport run and a controlled distributed evaluation are the
next evidence needed. This report records a preliminary implementation baseline.

\section{Simulation evidence after version 0.1.0 (31 October 2026)}
\label{sec:simulation-update}
This section summarizes the repository's dated validation record, rather than
replacing the historical experiment above. Its source snapshot is commit
\code{65da1ea} \cite{simulation_record}.
The virtual-time sweep and live ArduPilot software-in-the-loop (SITL) traffic
are distinct evidence types. Here, ``live'' means execution of simulation
processes on one host; it does not mean a physical vehicle or wireless link.
The numerical results below are transcribed from that record, not newly
measured or independently re-scored for this manuscript update.

\subsection{Receive-only mapping and declared ground truth}
The subsequent UDP MAVLink bridge maps \code{GLOBAL\_POSITION\_INT} position,
heading and horizontal speed, together with the latest
\code{SYS\_STATUS} battery percentage, to schema~v1. It withholds samples
when the battery is unknown or too old, heading is unknown, the source ID
does not match, or converted values violate the schema. Framing and mapping
counters distinguish these cases from monitor rejections. The bridge has
no encoder or send call; an automated test enforces that boundary. MAVProxy's
stream requests are external to the measured receive-only bridge.
Mapped sequence numbers follow arrival order, not the vehicle's source order.
Consequently this path cannot evaluate rejection of reordered source samples.

The seeded virtual-time impairment model declares fixed delay $\delta$,
uniform added jitter in $[0,j]$, independent loss, Gilbert--Elliott burst loss,
and outage windows $[u,v)$. Its documented matching rule associates an outage
with a STALE episode beginning in
\begin{equation}
u\leq t_{\mathrm{STALE}}\leq v+\tau+\delta+j+h.
\label{eq:outage-match}
\end{equation}
For example, with $[u,v)=[20,24)$ s, $\tau=1$ s, $\delta=0.1$ s,
$j=0.3$ s and $h=0.1$ s, the matching interval is $[20,25.5]$ s.
This is an experiment-scoring convention, not a physical fault diagnosis.
Detection latency is measured from outage onset $u$; recovery latency is
measured from outage end $v$. An outage beginning while already STALE is
reported separately, without a new detection latency. Unmatched STALE
episodes give the unexplained-STALE rate; in this model these are the
false-alarm measure relative to declared outages.

\subsection{Synthetic threshold sweep}
Table~\ref{tab:synthetic-sweep} condenses the documented sweep: 20 seeds per
row, 300 s per run, 10 Hz source, $h=0.05$ s, and two 3 s outages plus one
0.8 s outage per run. The default strict sequence-order policy was used.
The complete sweep and scenario-generation method remain in the pinned
repository record \cite{simulation_record}.

\begin{table}[htbp]
\centering\small
\setlength{\tabcolsep}{4pt}
\begin{tabular}{@{}p{35mm}rrrrrr@{}}
\toprule
Synthetic profile & $\tau$ (s) & 3 s: new+prior & 0.8 s & $\bar D$ (s) & $U$/h & Rej. \\
\midrule
Clean & 0.5 & $40+0$ & 20/20 & 0.40 & 0 & 0\% \\
Clean & 1.0 & $40+0$ & 0/20 & 0.90 & 0 & 0\% \\
Delay 0.1 s, jitter 0.3 s & 1.0 & $40+0$ & 3/20 & 1.19 & 0 & 25\% \\
Independent loss 30\% & 1.0 & $40+0$ & 3/20 & 0.86 & 0 & 0\% \\
Bursts, mean 1 s & 1.0 & $37+3$ & 6/20 & 0.83 & 240 & 0\% \\
Bursts, mean 1 s & 2.0 & $38+2$ & 0/20 & 1.78 & 71 & 0\% \\
\bottomrule
\end{tabular}
\caption{Reported synthetic sweep, not a network measurement.
Of 40 long outages, ``new+prior'' separates newly detected from already-STALE
outages; none are missed in these selected rows. $\bar D$ is mean detection
latency from outage onset for newly detected long outages, $U$/h is
unexplained STALE per hour, and Rej. is the out-of-order rejection percentage.
Short-outage misses are the complement of the 0.8 s column.}
\label{tab:synthetic-sweep}
\end{table}

In these scenarios, extending the timeout trades longer detection and more
missed short outages for fewer burst-induced STALE episodes. Average loss
alone does not characterize burst behavior. Delay and jitter increased
mean onset-referenced detection from 0.90 to 1.19 s in the selected rows;
this observed change is not a universal additive guarantee. Strict sequence
ordering rejected about a quarter of jittered arrivals in that row. A later
opt-in \code{late\_samples=liveness} policy was evaluated separately; these
numbers must not be attributed to that policy or to the MAVLink bridge.

\subsection{Clean-outage SITL observations and the conditional bound}
The 27 October record describes stationary and moving ArduPilot
\code{Rover-4.6.3} runs through MAVProxy~1.8.75 (pymavlink~2.4.50), a
one-way outage relay and the bridge, all on loopback. ArduPilot was pinned
to \code{3fc7011a7d3dc047cbb17d8bd98ee94577d144c6}. The bridge used
$\tau=1.0$ s and a nominal evaluation period of 0.5 s. Four declared
outages were matched to four STALE transitions; replay reproduced the
recorded state sequences with maximum transition differences of 0.24 ms
(stationary) and 1.1 ms (moving). Appendix~\ref{app:simulation-tables}
retains the four individual observations.

Using the last accepted arrival $r_\ast$ as the reference event, the reported
interval $t_{\mathrm{detect}}-r_\ast$ was 1.206--1.257 s. Substitution of
$\tau=1.0$ s and the assumed $h=0.5$ s in Equation~\eqref{eq:bound} gives
$[1.0,1.5]$ s, which contains those observations. This is consistency in
four observed outages, not a measured worst-case scheduler bound. A nominal
timer period does not establish the maximum evaluation gap required by
Equation~\eqref{eq:bound}. The onset-referenced detection intervals were
instead 1.093--1.104 s; the two reference events must not be interchanged.

Battery gating delayed recovery after a 5 s outage: a new position preceded
a fresh battery report. Initial $(0,0)$ positions were accepted before the
EKF origin existed, demonstrating that schema validity does not imply a
valid position fix. The initial 1 Hz position stream also straddled the
then-default 1 s threshold. The current bridge default is 3 s; historical
results at 1 s are not results at that new default.

\subsection{Reported live-impairment comparison and audit limits}
Issue~\#33 is closed and its implementation, PR~\#38, is merged; the
outline's ``results pending'' status is historical. The validation record
reports 20 SITL runs: four profiles, five seeds each, 300 s per run,
4 Hz position traffic, $\tau=1.0$ s and a 0.1 s bridge tick. After a
30 s warm-up, position sample $k$ receives the same seeded fate as virtual
sample $k$. Other datagrams share declared outages and burst state, but
independent loss and jitter are applied to position samples only. This
modeling choice does not emulate arbitrary per-datagram network loss.

The comparison distinguishes live observations, a model driven by observed
source timing, and a nominal periodic-source model. Table~\ref{tab:live-model}
keeps that distinction. The record reports 23,975/23,999 nominal fate matches,
with differences at outage boundaries. With observed source timing, reported
recovery differences were within 3.5 ms and detection differences within
0.087 s. The record reports reproduction of 39/39 loss and 59/59 burst
unexplained-STALE events. In contrast, the nominal model predicted all five
short burst-profile outages detected, against three observed. Source phase,
rate offset and interval variation therefore matter within this setup.

\textbf{These aggregates are not independently auditable from the repository
alone.} The original 20 relay logs and bridge recordings are not committed,
and this update establishes no public recording archive or DOI. The
30 October comparison-tool fix rejects inconsistent or truncated relay
logs; its regression tests did not reproduce the historical SITL runs.
A final summary establishes internal log consistency, not complete scenario
coverage. Neither the merged PR nor green software tests substitute for
re-scoring the original experiment artifacts.

Before treating the comparison as independently reproduced, recover and
archive each original scenario (including seed), complete relay log with
its terminal summary, bridge recording, JSON comparison report, exact
software versions and hashes. Verify scenario identity, log consistency,
scenario coverage and exit conditions with the current tools. Do not
reconstruct missing summaries. If original artifacts are unavailable,
retain the ``reported, not independently re-scored'' qualification and
perform separately dated replacement runs rather than presenting them
as the historical recordings.

\subsection{Limits and publication boundary}
These experiments cover one simulated vehicle on one host, loopback delivery
through MAVProxy, and runs of at most 300 s. The impairment comparison covers
one source rate and one timeout, with five seeds per profile; its rates
retain sampling uncertainty. It does not validate radio propagation, a real
network, a companion computer, physical vehicle dynamics, PX4, signed
MAVLink, or multi-vehicle SITL. Physical-link and field validation remain
future work. Arrival recency establishes neither measurement capture age
nor navigation accuracy, vehicle safety, autonomy, or research novelty.

Completion of \#33 permits discussion of reported simulation results, not
broader claims. Submission of a simulation-only report need not wait for
physical results, but remains subject to author review, the disclosed
artifact audit gap, and the author's publication and licence decisions.
No arXiv submission or prior arXiv version is established by this update.

\section*{Idea attribution}
The initial project idea is credited to \textbf{Christos Paparistodemou}, following
the project initiator's attribution. His public ResearchGate profile is
\url{https://www.researchgate.net/profile/Christos-Paparistodemou}.
The link identifies the contributor; the idea credit is supplied by the project
initiator. Project development and this report are presented under OpenPELAOS
by Charalampos Nadiotis.

\section*{Reproducibility and artifact availability}
The software targets Python~3.10 or newer; the authoring run used Python~3.12.14.
The optional transport target is ROS~2 Jazzy on Ubuntu~24.04. For the current repository layout, first run\\
\code{cd robotics/ros3\_limassol}; then:
\begin{quote}\small
\code{python -m ros3\_limassol.demo}\\
\code{python -m ros3\_limassol.demo --json}\\
\code{python -m unittest discover -s tests -v}
\end{quote}
The accompanying implementation archive includes the Apache-2.0 license,
contributor instructions, tests, and validation record. The public repository is
\url{https://github.com/n0711/ROS3-Limassol}; the robotics package is under
\code{robotics/ros3\_limassol}. The editable XeLaTeX manuscript includes its
bibliography. No arXiv identifier has been assigned at preparation time.

\textbf{Publication-preparation note (25 October 2026).}
Sections 1--7 retain the original 23 October 2026 version-0.1.0 experiment and
validation record. Subsequent software development reached version~0.3.1 at
commit \code{272a3b09c2271c56325683f8818e5436a377752f}, adding recording/replay,
a live receive-only recorder, and Humble compatibility work. A local rerun of
that commit discovered 45 tests: 42 passed and three ROS integration tests were
skipped because \code{rclpy} was unavailable. These counts must not be confused
with the original 18-test baseline. The repository validation record documents
CI evidence separately. This preparation corrects artifact availability and
font portability; it does not establish field performance or new novelty.

\clearpage
\appendix
\section{Reported SITL observations}
\label{app:simulation-tables}
These tables transcribe the 27 October validation record at the pinned
source snapshot \cite{simulation_record}; they are not new measurements.

\begin{table}[htbp]
\centering\small
\begin{tabular}{@{}lrrrr@{}}
\toprule
Run & Outage (s) & $D_{\mathrm{arrival}}$ (s) & $D_{\mathrm{onset}}$ (s) & $R$ (s) \\
\midrule
Stationary & 30--35 & 1.257 & 1.103 & 0.324 \\
Stationary & 70--72 & 1.206 & 1.104 & 0.083 \\
Moving & 30--35 & 1.239 & 1.093 & 0.324 \\
Moving & 70--72 & 1.232 & 1.093 & 0.085 \\
\bottomrule
\end{tabular}
\caption{Clean-outage SITL: $D_{\mathrm{arrival}}$ is time since the last
accepted arrival; $D_{\mathrm{onset}}$ is time since scheduled outage onset;
$R$ is time from traffic resumption to HEALTHY. All four outages were detected.}
\label{tab:clean-sitl}
\end{table}

\begin{table}[htbp]
\centering\small
\setlength{\tabcolsep}{4pt}
\begin{tabular}{@{}p{28mm}rrrr@{}}
\toprule
Profile & $\bar D$ (s) & $\bar R$ (s) & Short detected & $U$/h \\
\midrule
Clean & 0.81/0.80/0.80 & 0.15/0.15/0.00 & 5/5/5 & 0/0/0 \\
Delay 0.1 s, jitter 0.3 s & 1.09/1.09/1.07 & 0.38/0.38/0.24 & 5/5/5 & 0/0/0 \\
Independent loss 30\% & 0.65/0.62/0.60 & 0.22/0.22/0.05 & 5/5/5 & 94/94/79 \\
Bursts, mean 1 s & 0.79/0.80/0.77 & 0.11/0.11/0.25 & 3/3/5 & 142/142/139 \\
\bottomrule
\end{tabular}
\caption{Reported 20-run comparison. Each triplet is live / model with
observed source timing / nominal periodic model, in that order.
$\bar D$ and $\bar R$ concern the two 3 s outages per run; detection is
onset-referenced. All three methods report 10/10 long outages per profile.
Short counts are out of five 0.8 s outages; $U$/h is unexplained STALE per hour.
The nominal model uses \code{late\_samples=liveness}; the bridge numbers
samples in arrival order. Raw per-run artifacts remain unavailable here.}
\label{tab:live-model}
\end{table}

\begin{thebibliography}{9}\small\raggedright
\urlstyle{same}
\setlength{\itemsep}{3pt}
\bibitem{artifact}
C.~Nadiotis. \emph{ROS3-Limassol: OpenPELAOS Development Starter}.
Version~0.1.0, source package and \code{docs/VALIDATION.md}, 23 October 2026.
Software artifact, Apache-2.0. Public repository:
\url{https://github.com/n0711/ROS3-Limassol}.
The original report describes version~0.1.0; subsequent development is identified
separately in the publication-preparation note. No archival DOI assigned.
\bibitem{macenski2022}
S.~Macenski, T.~Foote, B.~Gerkey, C.~Lalancette, and W.~Woodall,
Robot Operating System 2: Design, architecture, and uses in the wild.
\emph{Science Robotics}, vol.~7, no.~66, article eabm6074, 2022.
\href{https://doi.org/10.1126/scirobotics.abm6074}{doi:10.1126/scirobotics.abm6074}.
Author manuscript: \url{https://arxiv.org/abs/2211.07752}.
\bibitem{ros_python}
Open Robotics. Writing a simple publisher and subscriber (Python).
\emph{ROS 2 Documentation: Jazzy}.
\url{https://docs.ros.org/en/jazzy/Tutorials/Beginner-Client-Libraries/Writing-A-Simple-Py-Publisher-And-Subscriber.html}.
Accessed 23 October 2026.
\bibitem{ros_qos}
Open Robotics. Quality of Service settings. \emph{ROS 2 Documentation: Jazzy}.
\url{https://docs.ros.org/en/jazzy/Concepts/Intermediate/About-Quality-of-Service-Settings.html}.
Accessed 23 October 2026.
\bibitem{chandra1996}
T.~D.~Chandra and S.~Toueg. Unreliable failure detectors for reliable distributed
systems. \emph{Journal of the ACM}, vol.~43, no.~2, pp.~225--267, 1996.
\href{https://doi.org/10.1145/226643.226647}{doi:10.1145/226643.226647}.
Author-maintained record: \url{https://www.cs.cornell.edu/info/people/sam/FDpapers.html}.
\bibitem{python_time}
Python Software Foundation. Time access and conversions:
\code{time.monotonic()}. \emph{Python 3.12 Documentation}.
\url{https://docs.python.org/3.12/library/time.html#time.monotonic}.
Accessed 23 October 2026.
\bibitem{rfc8259}
T.~Bray, Ed. \emph{The JavaScript Object Notation (JSON) Data Interchange Format}.
RFC~8259, December 2017.
\url{https://www.rfc-editor.org/rfc/rfc8259.html}.
\bibitem{simulation_record}
C.~Nadiotis. \emph{ROS3-Limassol repository validation and impairment records}.
Snapshot \code{65da1ea17d2a1de68e24a8f6e3120bbfd2fd1e11}, inspected
31 October 2026. Files: \code{robotics/ros3\_limassol/docs/VALIDATION.md},
\code{IMPAIRMENT.md}, \code{MAVLINK.md} and \code{SITL.md}; issue~\#33,
merged PR~\#38 and paper outline~\#42. Reported simulation results;
no independently re-scored raw-run archive established in this draft.
\url{https://github.com/n0711/ROS3-Limassol/tree/65da1ea17d2a1de68e24a8f6e3120bbfd2fd1e11/robotics/ros3_limassol/docs}.
\end{thebibliography}
\end{document}
